\documentclass[12pt]{article}

% ===== Core math =====
\usepackage{amsmath, amssymb, amsthm}
\usepackage{mathtools}
% ===== Diagrams =====
\usepackage{tikz-cd}
\usepackage{tikz}
\usetikzlibrary{arrows.meta, positioning, calc, fit, backgrounds}
% ===== References =====
\usepackage{hyperref}
\usepackage{cleveref}
% ===== Graphics / layout =====
\usepackage{graphicx}
\usepackage[margin=1in]{geometry}
% ===== GrokRxiv sidebar =====
% NOTE: `everypage` is retained deliberately: it is part of the mandated
% package set for this series' GrokRxiv template, and its `\AddEverypageHook`
% gives a robust page-1 test for the DOI sidebar below. Modern LaTeX kernels
% emit a benign "legacy status" informational notice on load (not an error);
% the document compiles cleanly. A future migration target is the kernel hook
% `\AddToHook{shipout/background}`.
\usepackage{everypage}
\usepackage{xcolor}
% ===== Misc =====
\usepackage{booktabs}
\usepackage{array}
\usepackage{longtable}
\usepackage{enumitem}
\usepackage{mdframed}

% Ragged-right fixed-width column type for wrapping tables (avoids overfull
% boxes and hyphenation overruns in the dictionary tables).
\newcolumntype{L}[1]{>{\raggedright\arraybackslash}p{#1}}

\hypersetup{
  colorlinks=true,
  linkcolor=blue!55!black,
  citecolor=green!45!black,
  urlcolor=blue!65!black,
  pdftitle={Motives, Periods, and Amplitudes: The Coalgebraic Anatomy of Physical Representation},
  pdfauthor={Matthew Long}
}

% ===== Theorem environments =====
\newtheorem{theorem}{Theorem}[section]
\newtheorem{proposition}[theorem]{Proposition}
\newtheorem{lemma}[theorem]{Lemma}
\newtheorem{corollary}[theorem]{Corollary}
\theoremstyle{definition}
\newtheorem{definition}[theorem]{Definition}
\newtheorem{example}[theorem]{Example}
\newtheorem{axiom}[theorem]{Axiom}
\theoremstyle{remark}
\newtheorem{remark}[theorem]{Remark}

% ===== Status-box environment (S/H/P epistemic labels) =====
\newmdenv[
  linecolor=gray!55,
  backgroundcolor=gray!7,
  linewidth=0.6pt,
  skipabove=8pt,
  skipbelow=8pt,
  innertopmargin=6pt,
  innerbottommargin=6pt
]{statusbox}

% ===== Notation macros =====
\providecommand{\llbracket}{[\![}
\providecommand{\rrbracket}{]\!]}
\newcommand{\Mot}{\mathrm{Mot}}
\newcommand{\Real}{\mathrm{Real}}
\newcommand{\Obs}{\mathrm{Obs}}
\newcommand{\per}{\mathrm{per}}
\newcommand{\Am}{\mathcal{A}}
\newcommand{\Amm}{\mathcal{A}^{\mathfrak{m}}}
\newcommand{\Lie}{\mathcal{L}}
\newcommand{\LG}{\mathcal{L}_{G}}
\newcommand{\Dom}{\mathsf{Dom}}
\newcommand{\Math}{\mathsf{Math}}
\newcommand{\Phys}{\mathsf{Phys}}
\newcommand{\Rep}{\mathsf{Rep}}
\newcommand{\Gpd}{\mathsf{Gpd}}
\newcommand{\Vect}{\mathsf{Vect}}
\newcommand{\Betti}{\mathrm{B}}
\newcommand{\dR}{\mathrm{dR}}
\newcommand{\Hdg}{\mathrm{Hdg}}
\newcommand{\et}{\text{\'et}}
\newcommand{\Li}{\mathrm{Li}}
\newcommand{\wt}{\mathrm{wt}}
\newcommand{\dep}{\mathrm{dep}}
\newcommand{\coA}{\Delta}
\newcommand{\cobr}{\delta}
\newcommand{\Prim}{\mathrm{Prim}}
\newcommand{\Zmot}{\mathcal{H}}
\newcommand{\Qbb}{\mathbb{Q}}
\newcommand{\Zbb}{\mathbb{Z}}
\newcommand{\Cbb}{\mathbb{C}}
\newcommand{\Rbb}{\mathbb{R}}
\newcommand{\Pbb}{\mathbb{P}}
\newcommand{\Gm}{\mathbb{G}_m}
\newcommand{\id}{\mathrm{id}}
\newcommand{\Hom}{\mathrm{Hom}}
\newcommand{\ol}[1]{\overline{#1}}
% Status labels: \ensuremath{\mathsf{...}} so they render identically in
% text and math mode without font-substitution warnings (reviewer, minor 3).
\newcommand{\Slab}{\ensuremath{\mathsf{S}}}
\newcommand{\Hlab}{\ensuremath{\mathsf{H}}}
\newcommand{\Plab}{\ensuremath{\mathsf{P}}}

% ===== GrokRxiv DOI sidebar =====
\definecolor{grokgray}{RGB}{110,110,110}
\AddEverypageHook{%
  \ifnum\value{page}=1
    \begin{tikzpicture}[remember picture, overlay]
      \node[
        rotate=90,
        anchor=south,
        font=\Large\sffamily\bfseries\color{grokgray},
        inner sep=0pt
      ] at ([xshift=38pt, yshift=0.52\paperheight]current page.south west)
      {GrokRxiv:2026.07.motives-periods-amplitudes\quad
       [\,math-ph\,]\quad
       04 Jul 2026};
    \end{tikzpicture}
  \fi
}

\title{\bfseries Motives, Periods, and Amplitudes:\\ The Coalgebraic Anatomy of Physical Representation\\[4pt]
\large Part II of \emph{A Math$\rightarrow$Physics Representation Library}}

\author{Matthew Long \\
\textit{The YonedaAI Collaboration} \\
\textit{YonedaAI Research Collective} \\
Chicago, IL \\
\texttt{research@yonedaai.com} $\cdot$ \url{https://yonedaai.com}}

\date{4 July 2026}

\begin{document}
\maketitle

\begin{abstract}
This is Part II of the modular series \emph{A Math$\rightarrow$Physics Representation Library}, whose governing thesis is that a large class of physical quantities are not merely \emph{described by} mathematics but are \emph{realized} from structured mathematical information along an explicit pipeline
$M \xmapsto{\Phi} \Phi(M) \xmapsto{\Real_\alpha} \text{phys.\ rep.} \xmapsto{\Obs} \text{meas.}$
Part I built the ambient \emph{representation stack} and left the realization channels $\Real_\alpha$ and the decomposition operation as abstract data. This paper supplies their first concrete instantiation. We formalize motives $\Mot(X)$ as the \emph{functional essence} of a mathematical object (a status-\Hlab{} semantic layer, carefully distinguished from Grothendieck's standard definition), we make the realization channels the named cohomological functors $\Real_\alpha$ for $\alpha \in \{\Betti,\dR,\Hdg,\et\}$, and we make the observable the period map $\per$. The central object is the \emph{motivic amplitude object} $\Amm$ with $\Am = \per(\Amm)$, equipped with a coaction $\coA(\Amm) = \sum_i \Amm_{i,1}\otimes \Amm_{i,2}$ and a cobracket $\cobr\colon \Lie \to \Lambda^2\Lie$. We prove: (T1) a \emph{Conditional Amplitude Decomposition} theorem stating that a realization channel, realized as a homomorphism of the motivic-Galois Hopf algebras, transports the coaction into a decomposition of the realized amplitude; (T2) that the weight-graded primitives cogenerate the Goncharov Lie coalgebra $\LG(F)$, with weight-one primitives identified with $F^\times\otimes\Qbb$; (T3) multiplicativity of the period pairing under direct sum and tensor product; (T4) a \emph{Non-Tate Obstruction} limitation theorem showing that elliptic and other non-Tate motivic lifts escape the polylogarithmic anatomy; and (T5) a conditional (status \Hlab/\Plab) compatibility statement for the cosmic Galois action. Throughout we retain the S/H/P epistemic-status discipline of Part I, and we exhibit the coalgebraic anatomy as the concrete arithmetic mechanism behind cuts, discontinuities, factorization, and residues. We accompany the theory with a compiling Haskell formalization of the coalgebra/coproduct structure and an idiomatic Lean sketch of the period/amplitude/coalgebra types.
\end{abstract}

\tableofcontents

% =====================================================================
\section{Introduction}
\label{sec:intro}

\subsection{The representation library and where Part II sits}

The series \emph{A Math$\rightarrow$Physics Representation Library} formalizes the idea that
mathematics functions as the \emph{syntax} of physical representation: physical
observables are obtained as realizations of structured mathematical information. The
program is deliberately \emph{modular}. Rather than positing a single universal functor
from all of mathematics to all of physics (which we argue in \Cref{sec:discussion}
is unlikely to be well defined), we build a ladder of domain-indexed modules. Each
module is defensible on its own, and each composes with its neighbours to add an
emergent structural property.

Part I introduced the ambient formalism: a category $\Dom$ of mathematical domains,
a category or higher category $\Math_d$ of structures in each domain, a category $\Phys$
of physical-representation targets, and the bracket notation $\llbracket M\rrbracket_{\Phys}$
for the physical representation assigned to $M$. Its central construct is the
\emph{realization pipeline}
\begin{equation}
\label{eq:pipeline}
\begin{aligned}
M \in \Math_d \;\xmapsto{\;\Phi\;}\;& \Phi(M)\ \text{(abstract info)}
\;\xmapsto{\;\Real_\alpha\;}\; \text{phys.\ rep.}
\;\xmapsto{\;\Obs\;}\; \text{measurement},\\[2pt]
& \Obs_\alpha(M) = \Obs\!\big(\Real_\alpha(\Phi(M))\big),
\end{aligned}
\end{equation}
together with a three-level epistemic-status system that every module must carry
(\Cref{subsec:status}). In Part I the maps $\Phi$, $\Real_\alpha$, $\Obs$ are abstract:
placeholders awaiting concrete content.

\emph{This paper provides that content in the motivic channel.} We take $\Phi$ to be
``pass to the motive'' $X\mapsto \Mot(X)$, we take $\Real_\alpha$ to be the classical
realization functors of a motive (Betti, de Rham, Hodge, \'etale), and we take $\Obs$ to
be the period map $\per$. Under this dictionary the compressed master slogan of the
series, $\text{observable}=\Real(\text{abstract structure})$, becomes the motivic
refinement
\begin{equation}
\label{eq:masterII}
\boxed{\;\Am = \per(\Amm),
\qquad
\coA(\Amm) = \sum_i \Amm_{i,1}\otimes \Amm_{i,2},
\qquad
\cobr\colon \Lie\to\Lambda^2\Lie.\;}
\end{equation}
Here $\Amm$ is a \emph{motivic amplitude object} (a pre-numerical object living in a
category or comodule of motivic periods), $\Am$ is the physical amplitude obtained by
taking periods, $\coA$ is the coaction that exposes the amplitude's internal anatomy, and
$\cobr$ is the cobracket that isolates its primitive antisymmetric pieces.

\subsection{The emergent property added by Part II}

In the language of the modular ladder, each level adds an emergent property that its
predecessor could only gesture at. Part I could speak of ``realization channels $\Real_\alpha$''
only in the abstract. Part II makes those channels \emph{concrete and numerically
checkable}: the Betti, de Rham, Hodge, and \'etale realizations are genuine functors with
explicit comparison isomorphisms, and periods are numbers one can compute. Part II also
endows observables with an \emph{internal coalgebraic anatomy}: the coaction and
cobracket decompose an amplitude into primitive informational pieces. Where Part I's
Decomposition Axiom asserted only that ``compositional observables admit a decomposition
operation,'' Part II identifies that operation with the Goncharov coproduct and its
cobracket, and proves structural theorems about it. This is the sense in which
\emph{coalgebras are the decomposition laws of physics}.

\subsection{The epistemic-status discipline}
\label{subsec:status}

Following Part I we tag every dictionary entry and every claim with one of three
epistemic-status labels, and we regard these labels as \emph{part of the formalism},
not as decoration.

\begin{statusbox}
\begin{center}
\begin{tabular}{@{}c L{0.42\textwidth} L{0.40\textwidth}@{}}
\toprule
\textbf{Label} & \textbf{Meaning} & \textbf{Canonical example} \\
\midrule
\Slab{} & standard use in mathematics or mathematical physics & Goncharov coproduct on polylogarithms \\
\Hlab{} & strong heuristic representation-dictionary entry & motive as ``functional essence'' \\
\Plab{} & speculative philosophical ontology & physical reality as motivic realization \\
\bottomrule
\end{tabular}
\end{center}
\end{statusbox}

\noindent
Composite labels such as \Slab/\Hlab{} and \Hlab/\Plab{} occur where an entry is standard as
pure mathematics but heuristic or speculative in its \emph{physical} reading. We preserve
these composite labels verbatim. The discipline forces us to state, for every amplitude
class, whether the motivic coaction is an established fact (status \Slab, e.g.\ Panzer--Schnetz
on $\phi^4$ periods \cite{PanzerSchnetz}) or a conjectural extension (status \Hlab/\Plab,
e.g.\ the cosmic Galois group \cite{BrownCosmic}).

\subsection{Contributions}

\begin{enumerate}[leftmargin=1.6em]
\item A precise formulation of the \emph{functional-essence} reading of a motive as a
status-\Hlab{} semantic layer over Grothendieck's standard notion, with a boxed caution
(\Cref{rem:not-functions}) that prevents the reading from being mistaken for the
mathematical definition.
\item An instantiation of Part I's realization pipeline \eqref{eq:pipeline} in the motivic
channel, identifying $\Phi,\Real_\alpha,\Obs$ with the motive, the cohomological
realizations, and the period map (\Cref{sec:motives}, \Cref{sec:periods}).
\item A self-contained proof of the \emph{Conditional Amplitude Decomposition} theorem
(\Cref{thm:cad}), stating that a realization channel (realized as a homomorphism of the
motivic-Galois Hopf algebras induced by a Tannakian realization functor) transports the
coaction on $\Amm$ to a decomposition of the realized amplitude, together with an explicit
statement of where the S/H/P labels bite.
\item A structure theorem for the Goncharov Lie coalgebra (\Cref{thm:goncharov-gen}): the
weight-graded primitives cogenerate $\LG(F)$, and weight-one primitives are $F^\times\otimes\Qbb$.
\item Multiplicativity of the period pairing (\Cref{thm:period-mult}) and a
\emph{Non-Tate Obstruction} limitation theorem (\Cref{thm:nontate}) giving a citable,
rigorous form of the series' Limitation~4.
\item A compiling Haskell formalization of period data, motivic amplitude objects,
coactions and cobrackets, and the Goncharov Lie coalgebra, together with an idiomatic Lean
type-signature sketch (\Cref{sec:formalization}).
\end{enumerate}

\subsection{Outline}
\Cref{sec:framework} recalls the Part~I framework and fixes motivic notation.
\Cref{sec:motives} treats motives as functional essence and the realization functors.
\Cref{sec:periods} develops periods and proves multiplicativity of the period pairing.
\Cref{sec:amplitude} introduces motivic amplitude objects, coalgebraic anatomy, and proves
the Conditional Amplitude Decomposition theorem.
\Cref{sec:goncharov} constructs the Goncharov Lie coalgebra and proves the generation
theorem.
\Cref{sec:cuts} gives the physics dictionary: coproducts, cuts, factorization, residues,
positive geometries, and the Connes--Kreimer Hopf algebra.
\Cref{sec:nontate} proves the Non-Tate Obstruction theorem and states the conditional cosmic
Galois compatibility.
\Cref{sec:results} collects the results.
\Cref{sec:formalization} presents the Haskell and Lean formalizations.
\Cref{sec:discussion} discusses limitations and the bridges to Parts I and III, and
\Cref{sec:conclusion} concludes.

% =====================================================================
\section{Mathematical Framework}
\label{sec:framework}

\subsection{Recollection of the Part I formalism}

We briefly recall the ambient objects; see Part I for details.

\begin{definition}[Representation entry, Part I]
\label{def:repentry}
A \emph{representation entry} is a quadruple $E = (M,P,\tau,\sigma)$ where $M\in\Math_d$ is a
mathematical structure, $P\in\Phys$ is a proposed physical representation, $\tau$ is a
semantic translation datum (a functor, natural transformation, equivalence class of models,
interpretation map, or weaker relation), and $\sigma\in\{\Slab,\Hlab,\Plab\}$ is the
epistemic-status label.
\end{definition}

\begin{definition}[Representation prestack / stack, Part I]
\label{def:prestack}
A \emph{representation prestack} is a pseudofunctor
$\Rep_{\Phys}\colon \Dom^{\mathrm{op}}\to\Gpd$ assigning to each domain $d$ a groupoid of
representation entries and to a refinement $u\colon e\to d$ a restriction functor
$u^\ast$. Given a Grothendieck topology on $\Dom$, $\Rep_{\Phys}$ is a
\emph{representation stack} if compatible families of local entries glue to global entries
uniquely up to coherent isomorphism.
\end{definition}

Part I imposes four axioms on the pipeline \eqref{eq:pipeline}, each of which Part II will
instantiate concretely:

\begin{axiom}[Realization]\label{ax:real}
$\text{observable} = \Obs\circ\Real_\alpha(\text{abstract structure})$.
\end{axiom}
\begin{axiom}[Equivalence]\label{ax:equiv}
Mathematical descriptions related by an equivalence in the automorphism/gauge groupoid of a
representation entry determine the same physical content at the chosen observational level.
\end{axiom}
\begin{axiom}[Locality/descent]\label{ax:local}
Physical representations must be locally assignable and globally coherent; failure of
descent is an obstruction, anomaly, or missing boundary datum.
\end{axiom}
\begin{axiom}[Decomposition]\label{ax:decomp}
Compositional, period-like observables admit a decomposition operation (boundary,
coproduct, coaction, spectral sequence, factorization) revealing lower-level informational
pieces.
\end{axiom}

The whole of Part II may be read as the statement: \emph{in the motivic channel,
\Cref{ax:decomp} is the Goncharov coproduct}, and \Cref{ax:real} is the period map.

\subsection{Motivic notation and standing conventions}

We work over a base field $k$ of characteristic zero, usually a number field, with a fixed
embedding $k\hookrightarrow\Cbb$. We write $\mathrm{Var}_k$ for smooth quasi-projective
varieties over $k$. We assume a Tannakian category $\mathrm{MM}(k)$ of mixed motives (or, in
the concrete cases we use, the category $\mathrm{MTM}(k)$ of mixed Tate motives, which is
unconditionally available over number fields by Levine, Deligne--Goncharov, and
Voevodsky, and whose period side is controlled by Brown \cite{BrownMTZ}). All tensor products
without a subscript are over $\Qbb$. For a graded vector space $V=\bigoplus_n V_n$ we write
$\wt(v)=n$ for $v\in V_n$ and call $n$ the \emph{weight}.

\begin{definition}[Realization functor]
\label{def:realization}
A \emph{realization} is an exact $\otimes$-functor $\Real_\alpha\colon \mathrm{MM}(k)\to
\mathcal{T}_\alpha$ to a Tannakian target $\mathcal{T}_\alpha$. The four classical
realizations are the \emph{Betti} realization $\Real_\Betti$ (singular cohomology,
target $\Qbb$-vector spaces with a Hodge/weight filtration), the \emph{de Rham} realization
$\Real_\dR$ (algebraic de Rham cohomology, target filtered $k$-vector spaces), the
\emph{Hodge} realization $\Real_\Hdg$ (mixed Hodge structures), and the \emph{\'etale}
realization $\Real_\et$ ($\ell$-adic Galois representations).
\end{definition}

The Betti and de Rham realizations of a motive are related by the \emph{comparison
isomorphism}
\begin{equation}
\label{eq:comparison}
\mathrm{comp}\colon \Real_\dR(M)\otimes_k\Cbb \;\xrightarrow{\ \sim\ }\;
\Real_\Betti(M)\otimes_\Qbb\Cbb,
\end{equation}
and \emph{periods} are the entries of the matrix of $\mathrm{comp}$ in chosen rational
bases. This is the arithmetic origin of \eqref{eq:masterII}: a period is what one
\emph{measures} when one realizes a motive both algebraically and topologically and compares
the two.

\begin{remark}[The motivic channel of the pipeline]
\Cref{def:realization} makes precise the sense in which Part II instantiates Part I's
$\Real_\alpha$: the abstract index $\alpha$ of \eqref{eq:pipeline} now ranges over the honest
set $\{\Betti,\dR,\Hdg,\et\}$, and each $\Real_\alpha$ is a genuine functor rather than a
placeholder. The remaining ``operational/quantum/thermodynamic'' channels of Part I are
treated in later parts of the series.
\end{remark}

% =====================================================================
\section{Motives as Functional Essence}
\label{sec:motives}

\subsection{The standard notion and the heuristic layer}

Grothendieck's motives are the universal cohomological invariants of algebraic varieties:
$\Mot(X)$ is designed so that every ``reasonable'' (Weil) cohomology theory factors through
a single realization functor out of the category of motives. Betti, de Rham, $\ell$-adic
\'etale, and crystalline cohomologies are then \emph{shadows} of one motivic object.

Our series overlays on this a semantic reading, which we flag as heuristic.

\begin{definition}[Functional-essence interpretation, status \Hlab]
\label{def:FE}
The \emph{functional-essence} interpretation assigns to a variety $X$ the semantic datum
\[
\mathrm{FE}(X) := \Mot(X),
\]
read as ``the invariant, pre-numerical structure from which the various functional readings,
cohomological realizations, periods, and numerical shadows of $X$ are realized.'' In the
pipeline \eqref{eq:pipeline} this is the concrete choice $\Phi(X)=\Mot(X)$.
\end{definition}

\begin{remark}[Caution: a motive is not literally ``the functions of'' $X$]
\label{rem:not-functions}
\begin{mdframed}[linecolor=red!45,linewidth=0.8pt,backgroundcolor=red!4]
It is essential not to misread \Cref{def:FE}. \emph{Mathematically}, a motive is not
``the functions of'' an object, and it is not the ring of functions $\mathcal{O}(X)$ or any
of its variants. A motive is a universal object unifying the different cohomology theories of
$X$; it is the source of $X$'s cohomological realizations, not its algebra of functions. The
reading ``functional essence'' is a \emph{proposed semantic layer} of this project (status
\Hlab), introduced to align motives with the pipeline's slogan
$\text{observable}=\Real(\text{abstract structure})$. Every use of the phrase ``functional
essence'' below inherits the status-\Hlab{} label of \Cref{def:FE} and must not be taken as a
mathematical identity.
\end{mdframed}
\end{remark}

\subsection{The realization chain}

Under \Cref{def:FE}, Part I's pipeline specializes to the \emph{motivic realization chain}
\begin{equation}
\label{eq:chain}
X \;\xmapsto{\ \Mot\ }\; \Mot(X) \;\xmapsto{\ \Real_\alpha\ }\; \Real_\alpha(\Mot(X))
\;\xmapsto{\ \Obs\ }\; \text{physical observable},
\end{equation}
with the identifications $\Phi=\Mot$, $\Obs=\per$. We display the chain, and its relation
to Part I's abstract pipeline, as a commuting diagram.

\begin{figure}[ht]
\centering
\begin{tikzcd}[column sep=large, row sep=large]
X \arrow[r, "\Mot"] \arrow[d, dashed, "\sim\ \text{(gauge/duality)}"'] &
\Mot(X) \arrow[r, "\Real_\alpha"] \arrow[d, "\text{same motive}"] &
\Real_\alpha(\Mot(X)) \arrow[r, "\per"] \arrow[d, "\text{comparison}"] &
\Am \arrow[d, equal] \\
X' \arrow[r, "\Mot"'] &
\Mot(X') \arrow[r, "\Real_{\alpha'}"'] &
\Real_{\alpha'}(\Mot(X')) \arrow[r, "\per"'] &
\Am
\end{tikzcd}
\caption{The motivic realization chain \eqref{eq:chain} and the Equivalence Axiom
(\Cref{ax:equiv}): if $X\sim X'$ under a motivic equivalence (e.g.\ a change of algebraic
model, a duality, or a gauge choice), the two rows realize the \emph{same} amplitude $\Am$.
Periods computed through different realization channels $\alpha,\alpha'$ agree via the
comparison isomorphism \eqref{eq:comparison}.}
\label{fig:chain}
\end{figure}

\begin{example}[The logarithm as a weight-two period]
\label{ex:log}
Let $X=\Gm=\Pbb^1\setminus\{0,\infty\}$ and consider the relative cohomology motive
$M=H^1(\Gm,\{1,x\})$ for $x\in k^\times$. Its de Rham side is spanned by the class of
$\tfrac{dt}{t}$ and its Betti side by the class of the path $\gamma$ from $1$ to $x$. The
period is
\[
\per\left(\int_\gamma \frac{dt}{t}\right) = \log x .
\]
This is the motivic incarnation of the statement ``$\log$ is a period''; the motive
$H^1(\Gm,\{1,x\})$ is an extension of $\Qbb(0)$ by $\Qbb(1)$, i.e.\ the Kummer motive. In
the functional-essence reading (status \Hlab), $\log x$ is a numerical shadow of the Kummer
motive, and the Kummer motive is the pre-numerical source common to all the ways one might
present $\log x$.
\end{example}

\begin{example}[Multiple polylogarithms]
\label{ex:polylog}
The multiple polylogarithms
\[
\Li_{n_1,\dots,n_d}(x_1,\dots,x_d) = \sum_{0<k_1<\cdots<k_d}
\frac{x_1^{k_1}\cdots x_d^{k_d}}{k_1^{n_1}\cdots k_d^{n_d}}
\]
are periods of the motivic fundamental group of $\Pbb^1\setminus\{0,1,\infty\}$ (and its
cyclotomic variants). Their weight is $n_1+\cdots+n_d$ and their depth is $d$. Multiple zeta
values $\zeta(n_1,\dots,n_d)=\Li_{n_1,\dots,n_d}(1,\dots,1)$ (with $n_d\ge 2$) are the
corresponding special values. These are the canonical periods populating amplitude
computations, and their motivic lifts generate the mixed Tate part of the coalgebra we build
in \Cref{sec:goncharov}.
\end{example}

% =====================================================================
\section{Periods and the Period Pairing}
\label{sec:periods}

\subsection{Period data and the period map}

\begin{definition}[Period datum]
\label{def:perioddatum}
A \emph{period datum} is a quadruple $\Pi = (X,D,\omega,\Gamma)$ where $X$ is a smooth
algebraic variety over $k$, $D\subset X$ is a normal-crossings divisor (the boundary),
$\omega$ is an algebraic differential form on $X\setminus D$, and $\Gamma$ is a relative cycle
in $H_\bullet(X(\Cbb),D(\Cbb);\Qbb)$ of complementary dimension. Its \emph{period} is
\begin{equation}
\per(\Pi) \;=\; \int_\Gamma \omega \;\in\;\Cbb.
\end{equation}
\end{definition}

Period data are the concrete carriers of the pairing between de Rham data ($\omega$) and
Betti data ($\Gamma$); the number $\per(\Pi)$ is exactly a matrix entry of the comparison
isomorphism \eqref{eq:comparison} for the motive $M=H^{\dim}(X,D)$. The
\emph{Kontsevich--Zagier} period ring $\mathcal{P}$ is the $\Qbb$-algebra generated by all
such numbers \cite{KontsevichZagier}; it contains $\ol{\Qbb}$, $\log 2$, $\pi$, $\zeta(3)$,
and the amplitude constants of perturbative quantum field theory.

\subsection{Multiplicativity of the period pairing}

The following theorem is the rigorous backbone that licenses treating $\Am=\per(\Amm)$ as
respecting the physical amplitude's additive and multiplicative structure.

\begin{theorem}[Multiplicativity of the period pairing]
\label{thm:period-mult}
Let $\Pi_1=(X_1,D_1,\omega_1,\Gamma_1)$ and $\Pi_2=(X_2,D_2,\omega_2,\Gamma_2)$ be period
data.
\begin{enumerate}[label=\textup{(\roman*)}, leftmargin=2.2em]
\item \textup{(Additivity.)} If $\Pi_1$ and $\Pi_2$ have the same $(X,D,\omega)$ and disjoint
cycles $\Gamma_1,\Gamma_2$, then the disjoint-union datum $\Pi_1\sqcup\Pi_2$ satisfies
$\per(\Pi_1\sqcup\Pi_2)=\per(\Pi_1)+\per(\Pi_2)$. More generally the period map is
$\Qbb$-bilinear in $(\omega,\Gamma)$.
\item \textup{(Multiplicativity.)} Define the \emph{product period datum}
\[
\Pi_1\otimes\Pi_2 := \big(X_1\times X_2,\ D_1\times X_2 \cup X_1\times D_2,\
\omega_1\wedge\omega_2,\ \Gamma_1\times\Gamma_2\big).
\]
Then
\[
\per(\Pi_1\otimes\Pi_2)=\per(\Pi_1)\cdot\per(\Pi_2).
\]
\end{enumerate}
\end{theorem}

\begin{remark}[What is and is not being asserted]
\label{rem:not-tautology}
The mathematical content of \Cref{thm:period-mult} is the pair of identities \textup{(i)} and
\textup{(ii)} --- in particular the multiplicativity \textup{(ii)}, which is the analytic
shadow of the K\"unneth decomposition and is not a formal consequence of the definitions. It
is these identities that make the evaluation map $\Pi\mapsto\per(\Pi)$ a \emph{ring}
homomorphism on the $\Qbb$-algebra of formal period data (modulo bilinearity, change of
variables, and Stokes). We emphasize that the \emph{surjectivity} onto $\mathcal{P}$ is
tautological, since $\mathcal{P}$ is \emph{defined} in \Cref{sec:periods} as the image of this
evaluation map; equivalently, \textup{(i)}--\textup{(ii)} are exactly what endows
$\mathcal{P}$ with its ring structure. The nontrivial statement is thus that the target is
closed under multiplication with the product realized by the product period datum, not that
the map hits everything.
\end{remark}

\begin{proof}
(i) The integral $\int_\Gamma\omega$ is $\Qbb$-linear in the chain $\Gamma$ and in the
form $\omega$ by definition of integration of forms over chains; for disjoint cycles
$\Gamma_1\sqcup\Gamma_2$ one has
$\int_{\Gamma_1\sqcup\Gamma_2}\omega=\int_{\Gamma_1}\omega+\int_{\Gamma_2}\omega$. This gives
additivity and, applied to both slots, $\Qbb$-bilinearity.

(ii) Write $p_j\colon X_1\times X_2\to X_j$ for the projections, so that
$\omega_1\wedge\omega_2 := p_1^\ast\omega_1\wedge p_2^\ast\omega_2$ is a form on
$(X_1\times X_2)\setminus(D_1\times X_2\cup X_1\times D_2)=(X_1\setminus D_1)\times
(X_2\setminus D_2)$. By the Fubini theorem for the product of the chains $\Gamma_1$ and
$\Gamma_2$ (both of complementary dimension in their factors, so that
$\Gamma_1\times\Gamma_2$ has complementary dimension in the product),
\[
\int_{\Gamma_1\times\Gamma_2}p_1^\ast\omega_1\wedge p_2^\ast\omega_2
=\left(\int_{\Gamma_1}\omega_1\right)\left(\int_{\Gamma_2}\omega_2\right).
\]
This is the analytic form of the K\"unneth decomposition
$H^\bullet(X_1\times X_2)\cong H^\bullet(X_1)\otimes H^\bullet(X_2)$ at the level of the
comparison isomorphism: the period matrix of a tensor product of motives is the Kronecker
product of the period matrices, and the displayed entry is the product of the corresponding
entries. This establishes the two displayed identities (i) and (ii); their ring-theoretic
reading is recorded in \Cref{rem:not-tautology}.
\end{proof}

\begin{corollary}[Periods respect factorized amplitudes]
\label{cor:factorize}
If a physical amplitude $\Am$ factorizes as $\Am=\Am_1\cdot\Am_2$ with each $\Am_i$
realized by a period datum $\Pi_i$, and if the joint amplitude is realized by the product
period datum $\Pi_1\otimes\Pi_2$, then the motivic identity $\Amm=\Amm_1\otimes\Amm_2$
realizes $\Am=\per(\Amm)$ compatibly with the factorization. Thus the master formula
$\Am=\per(\Amm)$ of \eqref{eq:masterII} is consistent with amplitude factorization at
tree level and, more generally, at any kinematic factorization channel where the geometry
degenerates to a product.
\end{corollary}

\begin{proof}
Immediate from \Cref{thm:period-mult}(ii): $\per(\Amm_1\otimes\Amm_2)=
\per(\Amm_1)\per(\Amm_2)=\Am_1\Am_2=\Am$.
\end{proof}

% =====================================================================
\section[Amplitude Objects and Coalgebraic Anatomy]{Motivic Amplitudes and Coalgebraic Anatomy}
\label{sec:amplitude}

\subsection{Coalgebras, comodules, and the coaction}

We recall the coalgebraic vocabulary in the form we use it.

\begin{definition}[Coalgebra, comodule]
\label{def:coalg}
A \emph{coalgebra} over $\Qbb$ is a $\Qbb$-vector space $C$ with a coproduct
$\coA\colon C\to C\otimes C$ and counit $\varepsilon\colon C\to\Qbb$ satisfying
coassociativity $(\coA\otimes\id)\coA=(\id\otimes\coA)\coA$ and counitality
$(\varepsilon\otimes\id)\coA=\id=(\id\otimes\varepsilon)\coA$. A \emph{right comodule} over
a Hopf algebra $H$ is a space $A$ with a coaction $\coA\colon A\to A\otimes H$ satisfying
$(\coA\otimes\id_H)\coA=(\id_A\otimes\Delta_H)\coA$ and
$(\id_A\otimes\varepsilon_H)\coA=\id_A$.
\end{definition}

\begin{definition}[Hopf algebra of motivic periods]
\label{def:hopfperiods}
Let $\Zmot=\bigoplus_{n\ge 0}\Zmot_n$ be a \emph{connected graded commutative Hopf algebra}
of motivic-Galois type: $\Zmot_0=\Qbb\cdot 1$ (connectedness), the grading is by weight, and
the coproduct $\coA\colon\Zmot\to\Zmot\otimes\Zmot$, product, unit $1$, counit
$\varepsilon\colon\Zmot\to\Qbb$ (projection onto $\Zmot_0$), and antipode make $\Zmot$ a Hopf
algebra. For mixed Tate motives over a number field, $\Zmot$ is the Hopf algebra of the
motivic Galois group, isomorphic as a graded vector space to the tensor coalgebra on
$\bigoplus_{n\ge 1}\mathrm{Ext}^1(\Qbb(0),\Qbb(n))$. Write $\bar\Zmot=\ker\varepsilon=
\bigoplus_{n\ge 1}\Zmot_n$ for the augmentation ideal.
\end{definition}

\begin{definition}[Motivic amplitude object]
\label{def:motamp}
A \emph{motivic amplitude object} is an element $\Amm\in\Zmot$ of the Hopf algebra of
\Cref{def:hopfperiods}. It is equipped with the weight and depth gradings
$\wt,\dep\colon\Zmot\to\Zbb_{\ge 0}$, together with the coproduct
\begin{equation}
\label{eq:coaction}
\coA\colon \Zmot\longrightarrow \Zmot\otimes \Zmot,\qquad
\coA(\Amm)=\sum_i \Amm_{i,1}\otimes \Amm_{i,2},
\end{equation}
and the period map $\per\colon\Zmot\to\Cbb$ with $\per(\Amm)=\Am$ the physical amplitude.
More generally, when a comodule algebra $A$ over $\Zmot$ is augmented with a chosen algebra
projection $A\twoheadrightarrow\Zmot$ (as holds for the mixed Tate periods), its elements are
also motivic amplitude objects via the coaction $A\to A\otimes\Zmot$.
\end{definition}

Because $\Zmot$ is connected and graded, the \emph{reduced coproduct}
\begin{equation}
\label{eq:reduced}
\coA'\colon \bar\Zmot\longrightarrow \bar\Zmot\otimes\bar\Zmot,\qquad
\coA'(x) := \coA(x) - x\otimes 1 - 1\otimes x \quad (x\in\bar\Zmot),
\end{equation}
is well defined and lands in $\bar\Zmot\otimes\bar\Zmot$: connectedness forces every term of
$\coA(x)$ other than $x\otimes 1$ and $1\otimes x$ to have both tensor factors of positive
weight. (This is the point at which a general external comodule would fail: there is no
canonical $1\otimes\id$ term for a bare comodule; the connected-graded-Hopf-algebra hypothesis
supplies it.) The reduced coproduct is the physically meaningful part; its symmetrization and
antisymmetrization give the coproduct and cobracket viewpoints.

\begin{definition}[Cobracket]
\label{def:cobracket}
On the associated graded Lie coalgebra of indecomposables
$\Lie=\bar\Zmot/\bar\Zmot^2=\bigoplus_{n\ge1}\Lie_n$ of $\Zmot$ (the ``Lie coalgebra of
primitives''), the \emph{cobracket} is the composite
\begin{equation}
\label{eq:cobracket}
\cobr\colon \Lie \xrightarrow{\ \coA'\ } \Lie\otimes\Lie
\xrightarrow{\ \pi\ } \Lambda^2\Lie,\qquad \pi(a\otimes b)=\tfrac12(a\wedge b),
\end{equation}
where $\pi$ is antisymmetrization. It is graded, $\cobr(\Lie_n)\subseteq\bigoplus_{p+q=n}
\Lie_p\wedge\Lie_q$, and satisfies the co-Jacobi identity dual to the Jacobi identity.
\end{definition}

\begin{definition}[Coalgebraic anatomy]
\label{def:anatomy}
Given a class of observable objects $\mathcal{O}$ equipped with a coalgebraic operation
$\coA$ or $\cobr$, the \emph{coalgebraic anatomy} of $O\in\mathcal{O}$ is the tree (or,
after antisymmetrization, the collection of Lie words) of lower-complexity pieces obtained
by iterating the decomposition. The \emph{primitives} are the leaves: elements $x$ with
$\cobr(x)=0$, equivalently $\coA'(x)=0$.
\end{definition}

\subsection{The Conditional Amplitude Decomposition theorem}

We now prove the central structural result: any realization channel that is monoidal
transports the coalgebraic anatomy of $\Amm$ to a genuine decomposition of the realized
amplitude. This is the theorem that makes \Cref{ax:decomp} of Part I concrete.

\begin{theorem}[Conditional Amplitude Decomposition]
\label{thm:cad}
Let $\Zmot$ be the connected graded Hopf algebra of motivic periods
\textup{(\Cref{def:hopfperiods})} with coproduct $\coA$, and let $\Zmot_\alpha$ be the
corresponding connected graded Hopf algebra on the realization side, with coproduct
$\coA_\alpha$. Let
\[
r_\alpha\colon \Zmot\longrightarrow \Zmot_\alpha
\]
be the homomorphism of graded Hopf algebras induced on the motivic-Galois Hopf algebras by a
Tannakian realization functor $\Real_\alpha\colon \mathrm{MM}(k)\to\mathcal{T}_\alpha$ of
\Cref{def:realization}; concretely $r_\alpha$ is the grading-preserving $\Qbb$-linear algebra
map that intertwines the coproducts, i.e.\ the square
\begin{equation}
\label{eq:compat}
\begin{tikzcd}[column sep=huge, row sep=large]
\Zmot \arrow[r, "\coA"] \arrow[d, "r_\alpha"'] &
\Zmot\otimes\Zmot \arrow[d, "r_\alpha\otimes r_\alpha"] \\
\Zmot_\alpha \arrow[r, "\coA_\alpha"'] &
\Zmot_\alpha\otimes\Zmot_\alpha
\end{tikzcd}
\end{equation}
commutes. Then for every motivic amplitude object $\Amm\in\Zmot$ the realized amplitude
$r_\alpha(\Amm)$ carries the decomposition
\[
\coA_\alpha\big(r_\alpha(\Amm)\big)=\sum_i r_\alpha(\Amm_{i,1})\otimes r_\alpha(\Amm_{i,2}),
\]
natural in $\Amm$ and compatible with the weight and depth gradings. Taking $\alpha$ to be
the numerical period realization $\per$ recovers the coaction principle on physical
amplitudes: the Goncharov coproduct on polylogarithms and Brown's coaction on Feynman periods.
\end{theorem}

\begin{proof}
Everything takes place at the level of graded vector spaces and $\Qbb$-linear maps; the
realization functor $\Real_\alpha$ enters only through the induced Hopf-algebra homomorphism
$r_\alpha$, so there is no application of a functor to a vector-space element. Evaluate the
commuting square \eqref{eq:compat} on $\Amm\in\Zmot$. Along the top-right path,
\[
(r_\alpha\otimes r_\alpha)\big(\coA(\Amm)\big)
=(r_\alpha\otimes r_\alpha)\Big(\sum_i\Amm_{i,1}\otimes\Amm_{i,2}\Big)
=\sum_i r_\alpha(\Amm_{i,1})\otimes r_\alpha(\Amm_{i,2}),
\]
using \eqref{eq:coaction} and $\Qbb$-bilinearity of $\otimes$. Along the left-bottom path,
$\coA_\alpha\big(r_\alpha(\Amm)\big)$. Commutativity of \eqref{eq:compat} equates the two,
giving the claimed identity. Existence of $r_\alpha$ as a graded Hopf-algebra homomorphism is
exactly the statement that a Tannakian realization functor induces a homomorphism of the
associated motivic-Galois Hopf algebras (functoriality of the Tannakian dual); the square
\eqref{eq:compat} is then the compatibility of that homomorphism with the coproducts, which
is part of the definition of a Hopf-algebra homomorphism. Naturality in $\Amm$ is the
linearity of $r_\alpha$; compatibility with gradings holds because $r_\alpha$ is
grading-preserving (weight and depth are motivic invariants transported by realization). For
$\alpha$ the period realization $r_\alpha=\per$, the displayed identity is precisely the
statement that the numerical coproduct of a period is computed by the Goncharov/Brown
coaction, which is the established content of \cite{Goncharov2005, BrownFeynman, PanzerSchnetz}.
\end{proof}

\begin{remark}[Where the S/H/P labels bite]
\label{rem:labels-bite}
\Cref{thm:cad} is, as pure algebra, a formal property of Hopf-algebra homomorphisms and
their compatibility with coproducts --- it is \emph{always} true once the hypotheses hold.
Its \emph{physical} content is entirely in the antecedent ``$\Amm\in\Zmot$'': whether a given
physical amplitude actually admits a motivic lift lying in the Hopf algebra $\Zmot$ of
motivic periods. For Feynman integrals evaluating to multiple zeta
values or (mixed Tate) polylogarithms this is established (status \Slab; e.g.\
\cite{PanzerSchnetz}). For amplitudes requiring elliptic or higher motivic structures the
comodule hypothesis fails over the Tate Hopf algebra (\Cref{thm:nontate}); the decomposition
is then only heuristic (status \Hlab) until one passes to the appropriate elliptic coalgebra.
Thus \Cref{thm:cad} is a status-\Slab{} theorem about a status-conditional hypothesis; the
conditionality is exactly the epistemic-honesty content of the S/H/P system.
\end{remark}

\subsection{Iterated coactions and the anatomy tree}

Iterating the reduced coaction produces the full anatomy. Write
$\coA'^{(2)}=(\coA'\otimes\id)\coA'$, etc. Coassociativity ensures the iterates are
well-defined up to the coalgebra axioms, and the maximal iteration terminates because weight
strictly decreases in each nontrivial tensor factor.

\begin{proposition}[Termination and primitivity of the anatomy]
\label{prop:termination}
For a motivic amplitude object $\Amm$ of weight $n$, the iterated reduced coaction terminates
after at most $n$ steps, and its terminal pieces are primitives (weight-one for the mixed
Tate polylogarithmic coalgebra). Consequently the coalgebraic anatomy of $\Amm$ is a finite
tree whose leaves are primitive periods.
\end{proposition}

\begin{proof}
The reduced coaction is graded and lowers weight: if $\wt(\Amm)=n$ then every term
$\Amm_{i,1}\otimes\Amm_{i,2}$ of $\coA'(\Amm)$ has $\wt(\Amm_{i,1})+\wt(\Amm_{i,2})=n$ with
both factors of weight $\ge 1$ (the weight-zero part is exactly the trivial terms subtracted
in $\coA'$). Hence each factor has weight $\le n-1$. By induction on weight, after at most
$n-1$ further steps every factor reaches weight one, at which point $\coA'$ vanishes (weight
one cannot split into two positive weights summing to one). Weight-one indecomposables are
the primitives; for the mixed Tate polylogarithmic coalgebra these are the logarithms
$F^\times\otimes\Qbb$ (\Cref{thm:goncharov-gen}). Finiteness of the tree follows since each
node has finitely many children (the coaction sum is finite) and depth is bounded by $n$.
\end{proof}

% =====================================================================
\section{The Goncharov Lie Coalgebra}
\label{sec:goncharov}

\subsection{Construction}

We now construct the central object of Source~2 of the knowledge base, the Goncharov
Lie coalgebra of a field, and prove its generation theorem.

\begin{definition}[Goncharov Lie coalgebra]
\label{def:GLC}
Let $F$ be a field. The \emph{Goncharov Lie coalgebra} $\LG(F)=\bigoplus_{n\ge 1}\LG(F)_n$ is
the graded Lie coalgebra of framed mixed Tate motives over $F$, generated by the (motivic)
polylogarithm classes $\Li_n^{\mathfrak{m}}$ and their multiple variants, graded by weight
$n$, and equipped with the cobracket
\begin{equation}
\label{eq:GLCcobracket}
\cobr\colon \LG(F)_n \longrightarrow \bigoplus_{p+q=n,\ p,q\ge 1}
\LG(F)_p\wedge\LG(F)_q,
\end{equation}
induced from the Goncharov coproduct on iterated integrals by antisymmetrization and
projection to indecomposables. The pair $(\LG(F),\cobr)$ is a graded Lie coalgebra: $\cobr$
is antisymmetric and satisfies the co-Jacobi identity.
\end{definition}

The cobracket is the shadow, on primitives, of Goncharov's explicit coproduct formula for
the iterated integral $I(a_0;a_1,\dots,a_N;a_{N+1})$:
\begin{multline}
\label{eq:goncharov-coproduct}
\coA\, I(a_0;a_1,\dots,a_N;a_{N+1})
=\sum_{\substack{0=i_0<i_1<\cdots\\ <i_k<i_{k+1}=N+1}}
\left(\prod_{p=0}^{k} I(a_{i_p};a_{i_p+1},\dots,a_{i_{p+1}-1};a_{i_{p+1}})\right)\\
{}\otimes I(a_0;a_{i_1},\dots,a_{i_k};a_{N+1}),
\end{multline}
whose combinatorics (sub-sequences of interpolation points) is the arithmetic
mechanism behind ``cuts'' in the physics dictionary of \Cref{sec:cuts}.

\subsection{The cogeneration theorem}

\begin{theorem}[Weight-graded primitives cogenerate $\LG(F)$]
\label{thm:goncharov-gen}
Let $F$ be a number field (so that the freeness properties of the motivic Galois Lie algebra
of \Cref{sec:framework} hold, by Borel's theorem and Deligne--Goncharov). Let
$\LG(F)=\bigoplus_{n\ge 1}\LG(F)_n$ with cobracket \eqref{eq:GLCcobracket}, and set
$P_n:=\ker\!\big(\cobr|_{\LG(F)_n}\big)$ (the weight-$n$ primitives). Then:
\begin{enumerate}[label=\textup{(\alph*)}, leftmargin=2.2em]
\item the primitives $\bigoplus_n P_n$ cogenerate $\LG(F)$: every element of weight $n$ is
recovered from iterated cobrackets of primitives via the coradical filtration;
\item there is a canonical isomorphism $P_1\cong F^\times\otimes\Qbb$, sending the weight-one
primitive attached to $x\in F^\times$ to $\{x\}\mapsto\log x$ under the period map;
\item the weight filtration $W_\bullet\LG(F)_n$ given by $W_m=\bigoplus_{j\le m}\LG(F)_j$ is
exhaustive, and the associated graded is spanned by iterated cobrackets of weight-one
elements.
\end{enumerate}
\end{theorem}

\begin{proof}
\emph{(b)} Weight-one framed mixed Tate motives over $F$ are extensions of $\Qbb(0)$ by
$\Qbb(1)$, classified by $\mathrm{Ext}^1_{\mathrm{MTM}(F)}(\Qbb(0),\Qbb(1))\cong
F^\times\otimes\Qbb$ (the Kummer motives of \Cref{ex:log}). Weight one cannot split as
$p+q$ with $p,q\ge 1$, so $\cobr|_{\LG(F)_1}=0$ and $P_1=\LG(F)_1\cong F^\times\otimes\Qbb$;
the period of the Kummer motive attached to $x$ is $\log x$, giving the stated map.

\emph{(c)} The weight filtration is exhaustive by construction ($\LG(F)$ is a direct sum of
its graded pieces). For the associated graded, note that $\cobr$ is graded and, by the
structure of the Goncharov coproduct \eqref{eq:goncharov-coproduct}, its leading term in the
depth filtration expresses a weight-$n$ generator modulo lower depth as a sum of cobrackets
of strictly lower weight. Dualizing: the graded dual of $(\LG(F),\cobr)$ is a graded Lie
\emph{algebra} (the motivic Galois Lie algebra), which for a number field $F$ is a free
graded Lie algebra on generators in each weight dual to the motivic $\mathrm{Ext}$-groups
$\mathrm{Ext}^1_{\mathrm{MTM}(F)}(\Qbb(0),\Qbb(n))$ (freeness by Borel's regulator theorem
and Deligne--Goncharov; over an arbitrary field freeness may fail, which is why we restrict
to number fields here); a free graded Lie algebra is generated by its degree-wise generators,
and the cobracket structure on the dual coalgebra is therefore cogenerated by the
corresponding primitives.

\emph{(a)} Combine (b) and (c) with the coradical filtration of a graded connected coalgebra:
for such a coalgebra the primitives $\Prim=\ker\coA'$ cogenerate, i.e.\ the intersection of
the kernels of all iterated reduced coactions is zero, so every element is detected by, and
reconstructed from, its iterated cobrackets landing in tensor powers of primitives. By
\Cref{prop:termination} the iteration terminates in weight $n$, and by (b) the terminal
primitives are $F^\times\otimes\Qbb$. Hence the weight-graded primitives cogenerate $\LG(F)$
as a graded Lie coalgebra.
\end{proof}

\begin{remark}[Physical reading, status \Hlab]
Under the amplitude dictionary, weight $=$ transcendentality $=$ loop/functional complexity,
and the cobracket $=$ single discontinuity/cut. \Cref{thm:goncharov-gen} then reads: \emph{all
amplitude complexity is generated by iterated single cuts of logarithmic (weight-one)
building blocks.} This is a candidate rigorous counterpart of the physics folklore that
multi-loop symbols are built from iterated discontinuities. We flag the physical reading as
status \Hlab; the mathematical statement (a)--(c) is status \Slab.
\end{remark}

\subsection{The symbol as maximal iteration}

The \emph{symbol} of a weight-$n$ polylogarithm is the maximal iteration of the cobracket
into weight-one pieces: $\mathcal{S}(\Amm)\in (F^\times\otimes\Qbb)^{\otimes n}$. It records
the leaves of the anatomy tree of \Cref{prop:termination} in order, and is computed by fully
un-shuffling the coproduct \eqref{eq:goncharov-coproduct}.

\begin{example}[Symbol of the dilogarithm]
\label{ex:dilog-symbol}
For the Bloch--Wigner motive of $\Li_2(x)$, the reduced coaction is
\[
\coA'\,\Li_2^{\mathfrak m}(x) = -\,\Li_1^{\mathfrak m}(x)\otimes\log^{\mathfrak m}(x)
= \log^{\mathfrak m}(1-x)\otimes\log^{\mathfrak m}(x),
\]
so the symbol is $\mathcal{S}(\Li_2(x))=-(1-x)\otimes x$. The two tensor slots are the two
weight-one primitives; the cobracket is their antisymmetrization
$\cobr\,\Li_2^{\mathfrak m}(x)=\log(1-x)\wedge\log(x)$. Physically the two entries are the two
branch loci $x=1$ and $x=0$ of the dilogarithm --- its cut structure --- exactly as the
dictionary of \Cref{sec:cuts} predicts.
\end{example}

% =====================================================================
\section{Coproducts, Cuts, and Factorization}
\label{sec:cuts}

\subsection{The decomposition dictionary}

The coalgebraic operations acquire direct physical readings. We reproduce the relevant part
of the operations dictionary with status labels.

\begin{statusbox}
\begin{center}
\small
\begin{tabular}{@{}L{0.34\textwidth} L{0.44\textwidth} c@{}}
\toprule
\textbf{Coalgebraic operation} & \textbf{Physical interpretation} & \textbf{Status} \\
\midrule
coproduct/coaction $\coA$ & full anatomy: how an amplitude splits into lower pieces & \Slab \\
cobracket $\cobr$ & antisymmetric primitive split; single cut/factorization channel & \Hlab \\
Goncharov coproduct \eqref{eq:goncharov-coproduct} & arithmetic anatomy of (multiple) polylogs & \Slab \\
symbol $\mathcal{S}$ & maximal discontinuity data; branch-cut structure & \Slab \\
primitive element & irreducible informational contribution & \Hlab \\
weight grading & transcendental complexity $=$ loop/functional depth & \Slab/\Hlab \\
depth grading & iterated-integral nesting $=$ nested propagation layers & \Slab/\Hlab \\
residue / pole & singular channel, threshold, on-shell condition & \Slab/\Hlab \\
monodromy & analytic continuation across a cut & \Slab \\
\bottomrule
\end{tabular}
\end{center}
\end{statusbox}

\subsection{Cuts and discontinuities}

The physical content of the cobracket is that a single cut of an amplitude (putting an
internal propagator on shell, or taking a discontinuity across a branch point) is
computed by the first entry of the coaction.

\begin{proposition}[Discontinuity computes the first coaction slot]
\label{prop:disc}
Let $\Amm$ be a motivic amplitude object with $\coA'(\Amm)=\sum_i\Amm_{i,1}\otimes\Amm_{i,2}$.
Under the period realization, the discontinuity of $\Am=\per(\Amm)$ across the branch point
associated to a weight-one primitive $\log s$ (where $s$ is a kinematic invariant vanishing
on the cut) is
\[
\mathrm{Disc}_s\,\Am = 2\pi i\sum_{i:\ \Amm_{i,2}=\log^{\mathfrak m} s}\per(\Amm_{i,1}),
\]
i.e.\ the first tensor factors paired with the second factor $\log s$. Iterating recovers
the full sequence of multiple discontinuities as deeper coaction slots.
\end{proposition}

\begin{proof}
The monodromy of $\per(\Amm)$ around $s=0$ acts on the Betti realization by the Picard--Lefschetz
transformation, whose logarithm is the weight-one nilpotent attached to $\log s$. Under the
motivic coaction the monodromy operator is dual to the second tensor slot: the coaction
$\coA'$ intertwines analytic continuation with the $H$-comodule structure, so the variation
$\Amm\mapsto\Amm(s e^{2\pi i})-\Amm(s)$ picks out exactly the terms whose second factor is
$\log^{\mathfrak m}s$, with coefficient $2\pi i$ from the period of $\Qbb(1)$
(namely $\per(\Qbb(1))=2\pi i$). Realizing gives the displayed formula; iterating on the
first factor $\Amm_{i,1}$ (itself a lower-weight amplitude object) yields the multiple
discontinuities as iterated first-slots.
\end{proof}

\subsection{Factorization and positive geometries}

At kinematic boundaries where an amplitude factorizes, the geometry underlying the period
degenerates. In the positive-geometry description \cite{ArkaniHamed}, an amplitude is the
canonical form of a positive geometry $(X,X_{\ge 0})$, and its boundary strata (the faces
of $X_{\ge 0}$) are themselves positive geometries whose canonical forms are the
factorized (residue) amplitudes.

\begin{example}[Boundary strata as coaction, status \Slab/\Hlab]
\label{ex:posgeom}
For the amplituhedron and associahedron picture of tree-level amplitudes, the residue of the
canonical form on a codimension-one face is the product of two lower-point canonical forms.
Matching this to \Cref{cor:factorize}: the codimension-one boundary is a product period datum
$\Pi_1\otimes\Pi_2$, so the residue equals $\per(\Amm_1)\per(\Amm_2)$. The coaction's leading
term $\coA'(\Amm)\supseteq\Amm_1\otimes(\cdots)$ organizes the poset of boundary strata
exactly as the coalgebra's coradical filtration organizes the anatomy tree. The identification
of ``boundary stratum'' with ``coaction slot'' is status \Slab{} where the geometry is a
genuine positive geometry with known canonical form and \Hlab{} where it is only conjectural.
\end{example}

\subsection{The Connes--Kreimer Hopf algebra}

On the renormalization side, the same coalgebraic grammar appears as the Connes--Kreimer Hopf
algebra $H_{\mathrm{CK}}$ of Feynman graphs \cite{ConnesKreimer}: the coproduct
\begin{equation}
\label{eq:CK}
\coA(\Gamma)=\Gamma\otimes 1 + 1\otimes\Gamma + \sum_{\gamma\subsetneq\Gamma}\gamma\otimes
\Gamma/\gamma
\end{equation}
sums over divergent subgraphs $\gamma$ (with $\Gamma/\gamma$ the contracted graph), and the
antipode $S$ implements the combinatorics of the renormalization $R$-operation
(BPHZ forest formula). This is the coalgebra of \emph{subdivergence bookkeeping}: where the
Goncharov coproduct decomposes the \emph{value} of an integral, the Connes--Kreimer coproduct
decomposes the \emph{graph}. The two are compatible via the period map, and their interaction
(the ``cosmic Galois'' action on renormalization) is developed in \Cref{sec:nontate}
and, in the Hopf-algebraic renormalization direction, in Part~VI of the series.

\begin{figure}[ht]
\centering
\begin{tikzcd}[column sep=huge, row sep=large]
\text{Feynman graph }\Gamma \arrow[r, "\coA_{\mathrm{CK}}"] \arrow[d, "\text{Feynman rules}=\per\circ\Mot"'] &
\displaystyle\sum \gamma\otimes\Gamma/\gamma \arrow[d, "\per\otimes\per"] \\
\text{amplitude }\Am=\per(\Amm) \arrow[r, "\coA_{\mathrm{motivic}}"'] &
\displaystyle\sum \per(\Amm_{i,1})\otimes\per(\Amm_{i,2})
\end{tikzcd}
\caption{Compatibility of the graph-level Connes--Kreimer coproduct \eqref{eq:CK} with the
value-level motivic coaction under the period map. The vertical arrows are the Feynman-rule
assignment of a period to a graph; the square commutes when the amplitude admits a motivic
lift (status \Slab{} for mixed Tate cases, e.g.\ \cite{PanzerSchnetz}).}
\label{fig:CK}
\end{figure}

% =====================================================================
\section{The Non-Tate Obstruction and the Cosmic Galois Group}
\label{sec:nontate}

\subsection{A limitation theorem}

Not every amplitude is polylogarithmic. The following theorem gives the series' Limitation~4
(``some amplitudes require elliptic, modular, non-Tate structures'') a precise, citable form.

\begin{theorem}[Non-Tate Obstruction]
\label{thm:nontate}
Let $\Amm$ be a motivic amplitude object whose motivic lift has, as a subquotient, the motive
$H^1(E)$ of an elliptic curve $E/k$ with $\mathrm{End}(E)=\Zbb$ (or, more generally, a
non-Tate simple motive). Then:
\begin{enumerate}[label=\textup{(\roman*)}, leftmargin=2.2em]
\item $\Amm$ is not a comodule element over the mixed Tate Hopf algebra $H_{\mathrm{MT}}$ of
iterated integrals of $d\log$-forms; equivalently, its coaction does not close within the
Tate coalgebra generated by $\{\log,\Li_n,\zeta\}$;
\item the weight-graded polylogarithmic anatomy of \Cref{thm:goncharov-gen} is unavailable:
the primitives are no longer exhausted by $F^\times\otimes\Qbb$, and the period matrix
contains the elliptic periods $\omega_1,\omega_2$ (and quasi-periods $\eta_1,\eta_2$) of $E$,
which are not $\Qbb$-linear combinations of multiple zeta values;
\item consequently one must replace $H_{\mathrm{MT}}$ by an elliptic-polylogarithm coalgebra
(iterated integrals on $E$, or on the modular curve), and \Cref{thm:cad} holds only relative
to that larger Hopf algebra.
\end{enumerate}
\end{theorem}

\begin{proof}
(ii) The category of mixed Tate motives over $k$ has simple objects only the Tate twists
$\Qbb(n)$; their periods lie in the $\Qbb$-algebra generated by $2\pi i$ and multiple zeta
(and polylog) values. The motive $H^1(E)$ is simple of rank two and \emph{not} isomorphic to
any $\Qbb(n)^{\oplus}$ (its Hodge numbers are $h^{1,0}=h^{0,1}=1$, whereas Tate motives are
of type $(p,p)$). Its period matrix is $\begin{psmallmatrix}\omega_1&\omega_2\\ \eta_1&\eta_2
\end{psmallmatrix}$ with $\det=2\pi i$ (Legendre relation); by transcendence results
(Chudnovsky, and the theory of periods of elliptic curves) $\omega_1$ is not a $\Qbb$-linear
combination of powers of $\pi$ and MZVs for $E$ without complex multiplication. Hence the
elliptic periods are genuinely new.

(i) If $\Amm$ were a comodule element over $H_{\mathrm{MT}}$, its coaction
$\coA(\Amm)\in\Amm\otimes H_{\mathrm{MT}}$ would express all its ``anatomy'' using only Tate
generators, forcing every subquotient of its motivic lift to be Tate --- contradicting the
presence of $H^1(E)$ as a subquotient. Concretely, the second tensor slots of $\coA(\Amm)$
would all be logarithms/zeta values, but the elliptic period $\omega_1$ appears in
$\per(\Amm)$ with a monodromy (an $\mathrm{SL}_2(\Zbb)$ modular transformation) that no Tate
comodule structure can reproduce.

(iii) The correct home is the Hopf algebra of iterated integrals of modular forms / functions
on $E$ (Brown, Levin--Racinet, Broedel--Duhr--Dulat--Tancredi). Over this larger algebra the
comodule hypothesis of \Cref{thm:cad} is restored, and the theorem applies verbatim with
$H_{\mathrm{MT}}$ replaced by $H_{\mathrm{ell}}$. This proves (iii).
\end{proof}

\begin{corollary}[Sharpness of the master formula]
The master formula $\Am=\per(\Amm)$ of \eqref{eq:masterII} remains valid for elliptic
amplitudes, but the \emph{coalgebraic anatomy} $\coA(\Amm)=\sum\Amm_{i,1}\otimes\Amm_{i,2}$
must be taken over $H_{\mathrm{ell}}$, not $H_{\mathrm{MT}}$. The status of the
polylogarithmic decomposition drops from \Slab{} to ``inapplicable'' precisely at the elliptic
threshold; this is the honest boundary of the mixed Tate story.
\end{corollary}

\subsection{The cosmic Galois group}

Brown's cosmic Galois group $G_c$ is a proposed pro-algebraic group acting on the ring of
(motivic) periods of a QFT, compatible with the coaction \cite{BrownCosmic, BrownPeriods}.
We record the conditional compatibility statement, carefully flagged.

\begin{proposition}[Cosmic Galois compatibility, status \Hlab/\Plab]
\label{prop:cosmic}
Suppose a cosmic Galois group $G_c$ acts on the Hopf algebra $H$ of motivic periods by Hopf
automorphisms, i.e.\ its coaction on the ring of motivic amplitude objects satisfies
$g\cdot\coA(\Amm)=(\mathrm{id}\otimes g)\,\coA(g\cdot\Amm)$ for $g\in G_c$. Then the action
descends through the period realization to an action on the numerical amplitude ring
commuting with the coaction-transported decomposition of \Cref{thm:cad}. In particular the
$G_c$-orbit of an amplitude constant (e.g.\ $\zeta(3)$) is closed under the coaction, and the
``small graph'' / weight-drop phenomena of $\phi^4$ periods are $G_c$-equivariant.
\end{proposition}

\begin{proof}[Proof (conditional/heuristic)]
Formally, if $G_c\subseteq\mathrm{Aut}_\otimes(\omega)$ is a subgroup of the Tannakian
automorphism group of the fibre functor $\omega=\per$ acting by Hopf automorphisms, then by
definition it commutes with the comodule structure $\coA$, and \Cref{thm:cad} transports this
equivariance through $\Real_\alpha=\per$. The content is entirely in the antecedent: a precise
definition of $G_c$ for a given QFT is open (it is conjecturally the motivic Galois group of
the category of periods appearing in that theory, cut down by physical constraints such as the
absence of $\zeta(2)$ / even weights in certain schemes). Because the antecedent is
conjectural, the whole statement carries status \Hlab/\Plab, matching the knowledge base's
labeling of the cosmic Galois entry. The verified instances --- the coaction on $\phi^4$
periods up to high loop order \cite{PanzerSchnetz} --- provide status-\Slab{} evidence for
special cases but do not establish the general $G_c$.
\end{proof}

\begin{remark}
\Cref{prop:cosmic} is deliberately the most speculative result in the paper. It is included
because the cosmic Galois group is the natural ``symmetry organizing amplitudes'' that the
whole coalgebraic picture points toward; but per the discipline of \Cref{subsec:status} we do
not present it as established. It sets up Part~III (variation of the geometry over moduli,
where $G_c$-equivariance becomes a constraint on Gauss--Manin connections) and Part~VI
(Hopf-algebraic renormalization, where $G_c$ acts on the Connes--Kreimer side).
\end{remark}

% =====================================================================
\section{Results}
\label{sec:results}

We collect the principal results of the paper and their epistemic status.

\begin{itemize}[leftmargin=1.6em]
\item \textbf{\Cref{thm:period-mult} (Multiplicativity, status \Slab).} The period pairing is
additive and multiplicative: $\per(\Pi_1\otimes\Pi_2)=\per(\Pi_1)\per(\Pi_2)$, so the period
map is a ring homomorphism onto the Kontsevich--Zagier period ring. This licenses
$\Am=\per(\Amm)$ to respect factorization (\Cref{cor:factorize}).
\item \textbf{\Cref{thm:cad} (Conditional Amplitude Decomposition, status \Slab{} on a
conditional hypothesis).} Any realization channel, as a graded Hopf-algebra homomorphism
$r_\alpha$ compatible with the coproduct,
transports $\coA(\Amm)=\sum_i\Amm_{i,1}\otimes\Amm_{i,2}$ to a decomposition of the realized
amplitude. This is the concrete form of Part I's Decomposition Axiom (\Cref{ax:decomp}). The
physical content is in whether a given amplitude admits a motivic lift in $\Zmot$
(\Cref{rem:labels-bite}).
\item \textbf{\Cref{prop:termination} (Termination, status \Slab).} The iterated coaction of a
weight-$n$ amplitude object terminates in at most $n$ steps at primitive leaves.
\item \textbf{\Cref{thm:goncharov-gen} (Cogeneration, status \Slab).} The weight-graded
primitives cogenerate the Goncharov Lie coalgebra $\LG(F)$, and weight-one primitives are
$F^\times\otimes\Qbb$ (logarithms). Physically (status \Hlab): all amplitude complexity is
built from iterated cuts of logarithms.
\item \textbf{\Cref{prop:disc} (Discontinuity, status \Slab).} The discontinuity of an
amplitude across a branch point is the first coaction slot paired with the corresponding
weight-one primitive.
\item \textbf{\Cref{thm:nontate} (Non-Tate Obstruction, status \Slab).} Elliptic (non-Tate)
motivic lifts escape the mixed Tate coalgebra; the polylogarithmic anatomy is unavailable and
one must pass to an elliptic-polylogarithm Hopf algebra. This is the rigorous form of the
series' Limitation~4.
\item \textbf{\Cref{prop:cosmic} (Cosmic Galois, status \Hlab/\Plab).} Conditional on a
precise definition of the cosmic Galois group, its action commutes with the coaction and
descends to the numerical amplitude ring. Presented as conjectural.
\end{itemize}

The logical dependency of the results is displayed below.

\begin{figure}[ht]
\centering
\resizebox{\textwidth}{!}{%
\begin{tikzcd}[ampersand replacement=\&, column sep=1.4em, row sep=1.5em]
\& \text{\Cref{def:coalg}: coalgebra/comodule} \arrow[d] \& \\
\text{\Cref{def:perioddatum}: period datum} \arrow[dr] \&
\text{\Cref{def:motamp}: motivic amplitude} \arrow[d] \arrow[dl] \&
\text{\Cref{def:GLC}: Goncharov coalgebra} \arrow[d] \\
\text{\Cref{thm:period-mult}} \arrow[d] \&
\text{\Cref{thm:cad}} \arrow[d] \&
\text{\Cref{thm:goncharov-gen}} \arrow[dl] \\
\text{\Cref{cor:factorize}} \&
\text{\Cref{prop:disc},\ \Cref{thm:nontate},\ \Cref{prop:cosmic}} \&
\end{tikzcd}}
\caption{Dependency graph of the paper's definitions, theorems, and propositions.}
\label{fig:deps}
\end{figure}

% =====================================================================
\section{Formalization: Haskell and Lean}
\label{sec:formalization}

\begin{sloppypar}
To honor the Curry--Howard--\allowbreak Lambek bridge of the series (propositions as types,
proofs as programs, realization as compilation), we accompany the mathematics with a compiling
Haskell formalization of the coalgebra/coproduct structure and an idiomatic Lean type sketch.
The full sources live in
\texttt{src/\allowbreak motives-\allowbreak periods-\allowbreak amplitudes/} and
\texttt{lean/\allowbreak motives-\allowbreak periods-\allowbreak amplitudes/}; we display the
load-bearing fragments.
\end{sloppypar}

\subsection{Haskell: the coalgebra and coproduct}

The Haskell package represents a graded coalgebra element as a formal $\Qbb$-linear
combination of \emph{words} in generators, with the Goncharov-style deconcatenation coproduct.
Weight is the word length (for the shuffle/deconcatenation model) and the reduced coproduct
strips the two trivial terms; primitivity is decidable.

\begin{verbatim}
-- Word in generators; weight = length; deconcatenation coproduct
newtype Word' g = Word' [g] deriving (Eq, Ord, Show)

-- Formal Q-linear combination (rationals as coefficients)
newtype LinComb g = LinComb (Map (Word' g) Rational)

-- Deconcatenation coproduct: Delta(w) = sum over splittings u | v
coproduct :: Ord g => Word' g -> [(Word' g, Word' g)]
coproduct (Word' xs) =
  [ (Word' (take k xs), Word' (drop k xs)) | k <- [0 .. length xs] ]

-- Reduced coproduct: drop the two trivial terms (empty | w) and (w | empty)
reducedCoproduct :: Ord g => Word' g -> [(Word' g, Word' g)]
reducedCoproduct w@(Word' xs) =
  [ p | p@(Word' a, Word' b) <- coproduct w, not (null a), not (null b) ]

-- Cobracket delta = antisymmetrization of the reduced coproduct.
-- NOTE: this returns a formal signed sum WITHOUT combining like terms;
-- algebraic simplification in the exterior algebra Lambda^2 L is omitted
-- (adequate for the demonstration, not a normal form).
cobracket :: Ord g => Word' g -> [(Word' g, Word' g, Rational)]
cobracket w =
  [ (a, b, 1) | (a, b) <- reducedCoproduct w ] ++
  [ (b, a, -1) | (a, b) <- reducedCoproduct w ]

-- Primitivity: delta(x) = 0  <=>  reduced coproduct empty  <=>  weight <= 1
isPrimitive :: Ord g => Word' g -> Bool
isPrimitive w = null (reducedCoproduct w)
\end{verbatim}

The demonstration program verifies, on the dilogarithm symbol of
\Cref{ex:dilog-symbol}, that (a) coassociativity holds numerically for the deconcatenation
coproduct, (b) weight-one words are exactly the primitives (\Cref{thm:goncharov-gen}(b)), and
(c) the iterated reduced coproduct of a weight-$n$ word terminates after $n-1$ steps at
weight-one leaves (\Cref{prop:termination}). It compiles with \texttt{ghc} and runs to print
the anatomy tree of a sample amplitude word.

\subsection{Lean: period, amplitude, and coalgebra types (best-effort sketch)}

The Lean file records the type signatures of the objects and the statement of
\Cref{thm:cad} as a definition-level goal (with \texttt{sorry} placeholders, as a
non-build-gated sketch).

\begin{verbatim}
structure PeriodDatum (X Form Cycle : Type) where
  space   : X
  divisor : X
  form    : Form
  cycle   : Cycle

-- per : PeriodDatum -> C  (integration pairing, abstracted)
def per {X F C : Type} (integrate : C -> F -> Complex)
    (Pi : PeriodDatum X F C) : Complex := integrate Pi.cycle Pi.form

-- A right H-comodule of motivic amplitude objects
structure MotivicAmplitude (H A : Type) where
  coaction : A -> List (A x H)     -- Delta(a) = sum a_{i,1} (x) a_{i,2}
  weight   : A -> Nat
  depth    : A -> Nat

-- Cobracket as antisymmetrized reduced coaction on primitives L
def cobracket {L : Type} (coact : L -> List (L x L)) (x : L) :
    List (L x L) := coact x  -- (antisymmetrization elided in sketch)

-- T1 (Conditional Amplitude Decomposition), signature-level.
-- The realization is a MAP between carrier types (a Hopf-algebra
-- homomorphism), NOT a functor applied to elements: this is exactly the
-- corrected formulation of Theorem 5.3.
theorem coaction_transports
    {H A R : Type} (ma : MotivicAmplitude H A)
    (Real : A -> R) (RealH : H -> R) (target : R -> List (R x R)) (a : A) :
    target (Real a)
      = (ma.coaction a).map (fun p => (Real p.1, RealH p.2)) := by
  sorry
\end{verbatim}

The Lean sketch is explicitly best-effort and not build-gated; it fixes the type-level
anatomy (period data, comodule coactions, cobrackets) so that a future Mathlib-backed
development can discharge the \texttt{sorry}s. Note that modelling the realization as an
honest map \texttt{Real : A -> R} between carrier types (rather than a functor applied to
elements) is precisely the type-level reflection of the Hopf-algebra-homomorphism formulation
of \Cref{thm:cad}; the \texttt{period\_multiplicative} statement in the same file is proved
outright (no \texttt{sorry}) from the Fubini/K\"unneth hypothesis.

% =====================================================================
\section{Discussion}
\label{sec:discussion}

\subsection{Limitations (reproduced and addressed)}

Per the series discipline, we restate the global limitations and indicate how Part II
addresses each.

\begin{enumerate}[leftmargin=1.6em]
\item \emph{Not every mathematical analogy is a physical theory.} The S/H/P labels are part of
the formalism (\Cref{subsec:status}); the strongest results here (\Cref{thm:period-mult},
\Cref{thm:cad}, \Cref{thm:goncharov-gen}) are status \Slab{} as mathematics, while their
physical readings are labelled \Hlab.
\item \emph{Motives are not literally ``the functions of'' an object.} Addressed head-on in the
boxed \Cref{rem:not-functions}; the functional-essence reading is status \Hlab{} throughout.
\item \emph{No single universal functor $\Math\to\Phys$.} Part II defines only the motivic
channel $\Real_\alpha$, $\alpha\in\{\Betti,\dR,\Hdg,\et\}$, over the domain of motivic period
theory, a local, domain-indexed construction rather than a global functor.
\item \emph{Not every observable is a period.} Made precise as the Non-Tate Obstruction
theorem (\Cref{thm:nontate}): elliptic and modular amplitudes require larger coalgebras.
\item \emph{Gauge redundancy $\ne$ physical symmetry.} Deferred to Part~VI, but foreshadowed in
\Cref{fig:chain}: the dashed equivalence $X\sim X'$ is a motivic (gauge/duality) equivalence
realizing the same amplitude.
\end{enumerate}

\subsection{Bridges to Parts I and III}

Part II \emph{builds on} Part I by instantiating its abstract pipeline: $\Phi=\Mot$,
$\Real_\alpha=$ the cohomological realizations, $\Obs=\per$, and \Cref{ax:decomp}$=$ the
Goncharov coproduct. It \emph{sets up} Part III by handing off the motivic amplitude objects
and their families: as kinematic parameters vary, the variety $X$ in a period datum
$\Pi=(X,D,\omega,\Gamma)$ moves in a moduli space, the period $\per(\Pi)$ becomes a
multivalued function on that base, and the coaction becomes a statement about the
Gauss--Manin/Picard--Fuchs connection governing the variation of Hodge structure. In the
ladder's language, Part III lets Part II's motives ``live over and vary across concrete
geometric moduli.''

\begin{figure}[ht]
\centering
\resizebox{\textwidth}{!}{%
\begin{tikzcd}[ampersand replacement=\&, column sep=1.8em, row sep=1.4em]
\textbf{Part I} \arrow[r, "\text{instantiate}"] \&
\textbf{Part II} \arrow[r, "\text{vary over moduli}"] \&
\textbf{Part III} \\
\Phi,\Real_\alpha,\Obs\ (\text{abstract}) \arrow[u, phantom] \&
\Mot,\{\Real_\Betti,\Real_\dR,\Real_\Hdg,\Real_\et\},\per \arrow[u, phantom] \&
\text{VHS, Gauss--Manin, Picard--Fuchs} \arrow[u, phantom] \\
\text{Decomposition Axiom} \arrow[u, phantom] \&
\text{Goncharov coproduct }\coA,\ \cobr \arrow[u, phantom] \&
\text{monodromy over moduli} \arrow[u, phantom]
\end{tikzcd}}
\caption{Position of Part II in the modular ladder: it makes Part I's abstract channels
concrete and hands amplitude objects to Part III as objects varying over moduli.}
\label{fig:ladder}
\end{figure}

\subsection{The four slogans}

The series' compressed theses read, in the light of Part II:
\emph{mathematics is the syntax of physical representation}; \emph{motives are functional
semantics} (\Cref{def:FE}, status \Hlab); \emph{periods are numerical measurements}
(\Cref{def:perioddatum}, \Cref{thm:period-mult}); \emph{coalgebras are decomposition laws}
(\Cref{thm:cad}, \Cref{thm:goncharov-gen}). Part II is precisely the module that turns the
last two slogans into theorems.

% =====================================================================
\section{Conclusion}
\label{sec:conclusion}

We have given the motivic channel of the \emph{Math$\rightarrow$Physics Representation Library}
its first concrete instantiation. Motives supply the pre-numerical functional essence
(status \Hlab, carefully distinguished from the mathematical definition); the four
cohomological realizations supply Part I's abstract channels $\Real_\alpha$; and the period
map supplies the observable. On this foundation the master formula $\Am=\per(\Amm)$ with its
coaction $\coA(\Amm)=\sum_i\Amm_{i,1}\otimes\Amm_{i,2}$ and cobracket
$\cobr\colon\Lie\to\Lambda^2\Lie$ becomes a working piece of mathematics: the period pairing is
multiplicative (\Cref{thm:period-mult}), the coaction transports through realizations
(\Cref{thm:cad}), the Goncharov Lie coalgebra is cogenerated by its weight-one logarithmic
primitives (\Cref{thm:goncharov-gen}), amplitude discontinuities are read off the first
coaction slot (\Cref{prop:disc}), and the whole polylogarithmic anatomy has a sharp boundary
at the elliptic threshold (\Cref{thm:nontate}). The most speculative extension, the cosmic
Galois symmetry (\Cref{prop:cosmic}), is flagged as such and handed forward.

The emergent property Part II adds to the ladder is the \emph{internal coalgebraic anatomy of
observables}: the observation that a physical amplitude carries a canonical decomposition
into primitive informational pieces, governed by a Lie coalgebra, rather than standing as an
indecomposable whole. In the next module this anatomy is set in motion over the moduli spaces
of algebraic geometry; in the capstone module the same Hopf-coalgebraic grammar returns as the
bookkeeping of gauge redundancy and renormalization. Across the modules, the same coalgebraic
operations take an observable apart into its primitive pieces.

% =====================================================================
\begin{thebibliography}{99}
\sloppy
\addcontentsline{toc}{section}{References}

\bibitem{KontsevichZagier}
M.~Kontsevich and D.~Zagier,
\emph{Periods}, in ``Mathematics Unlimited --- 2001 and Beyond,'' Springer, 2001, pp.~771--808.

\bibitem{Goncharov2001}
A.~B.~Goncharov,
\emph{Multiple polylogarithms and mixed Tate motives}, arXiv:math/0103059 (2001).

\bibitem{Goncharov2005}
A.~B.~Goncharov,
\emph{Galois symmetries of fundamental groupoids and noncommutative geometry},
Duke Math.\ J.\ \textbf{128} (2005), no.~2, 209--284; arXiv:math/0208144.

\bibitem{GSVV}
A.~B.~Goncharov, M.~Spradlin, C.~Vergu, and A.~Volovich,
\emph{Classical polylogarithms for amplitudes and Wilson loops},
Phys.\ Rev.\ Lett.\ \textbf{105} (2010), 151605; arXiv:1006.5703.

\bibitem{BEK}
S.~Bloch, H.~Esnault, and D.~Kreimer,
\emph{On motives associated to graph polynomials},
Comm.\ Math.\ Phys.\ \textbf{267} (2006), 181--225; arXiv:math/0510011.

\bibitem{BrownMTZ}
F.~Brown,
\emph{Mixed Tate motives over $\Zbb$},
Annals of Mathematics \textbf{175} (2012), no.~2, 949--976; arXiv:1102.1312.

\bibitem{BrownFeynman}
F.~Brown,
\emph{On the periods of some Feynman integrals}, arXiv:0910.0114 (2009).

\bibitem{BrownCosmic}
F.~Brown,
\emph{Feynman amplitudes, coaction principle, and cosmic Galois group},
Commun.\ Number Theory Phys.\ \textbf{11} (2017), 453--556; arXiv:1512.06409.

\bibitem{BrownPeriods}
F.~Brown,
\emph{Periods and Feynman amplitudes}, arXiv:1512.09265 (2015).

\bibitem{PanzerSchnetz}
E.~Panzer and O.~Schnetz,
\emph{The Galois coaction on $\phi^4$ periods},
Commun.\ Number Theory Phys.\ \textbf{11} (2017), 657--705; arXiv:1603.04289.

\bibitem{ConnesKreimer}
A.~Connes and D.~Kreimer,
\emph{Renormalization in quantum field theory and the Riemann--Hilbert problem I: the Hopf
algebra structure of graphs and the main theorem},
Comm.\ Math.\ Phys.\ \textbf{210} (2000), 249--273; arXiv:hep-th/9912092.

\bibitem{ArkaniHamed}
N.~Arkani-Hamed, Y.~Bai, and T.~Lam,
\emph{Positive geometries and canonical forms},
JHEP \textbf{11} (2017) 039; arXiv:1703.04541.

\bibitem{Deligne}
P.~Deligne and A.~B.~Goncharov,
\emph{Groupes fondamentaux motiviques de Tate mixte},
Ann.\ Sci.\ \'Ecole Norm.\ Sup.\ \textbf{38} (2005), 1--56.

\bibitem{BrownDepth}
F.~Brown,
\emph{On the decomposition of motivic multiple zeta values},
in ``Galois--Teich\-m\"uller theory and arithmetic geometry,''
Adv.\ Stud.\ Pure Math.\ \textbf{63} (2012), 31--58; arXiv:1102.1310.

\bibitem{BroedelEtal}
J.~Broedel, C.~Duhr, F.~Dulat, and L.~Tancredi,
\emph{Elliptic polylogarithms and iterated integrals on elliptic curves},
JHEP \textbf{05} (2018) 093; arXiv:1712.07089.

\bibitem{DuhrReview}
C.~Duhr,
\emph{Mathematical aspects of scattering amplitudes},
in ``Journeys Through the Precision Frontier: Amplitudes for Colliders'' (TASI 2014),
World Scientific, 2016, pp.~419--476; arXiv:1411.7538.

\end{thebibliography}

% =====================================================================
\appendix
\section[The Motivic and Amplitude Library]{The Motivic and Amplitude Library}
\label{app:library}

For reference we reproduce the motivic/amplitude dictionary of the series, with epistemic
status labels retained verbatim. This table is the Appendix~B of the source library,
restricted to Part II's domain.

\renewcommand{\arraystretch}{1.15}
\setlength{\tabcolsep}{4pt}
\small
\begin{longtable}{@{}L{0.26\textwidth} L{0.31\textwidth} L{0.29\textwidth} c@{}}
\toprule
\textbf{Mathematics} & \textbf{Physical representation} & \textbf{Decomposition meaning} & \textbf{Status} \\
\midrule
\endfirsthead
\toprule
\textbf{Mathematics} & \textbf{Physical representation} & \textbf{Decomposition meaning} & \textbf{Status} \\
\midrule
\endhead
\bottomrule
\endfoot
Field $F$ & Kinematic/coordinate domain & Allowed parameter values & \Slab/\Hlab \\
$F^\times$ & Nonzero scales or ratios & Energies, masses, cross-ratios & \Slab/\Hlab \\
Cross-ratio & Conformal invariant & Dimensionless observable & \Slab \\
Algebraic variety & Constraint solution space & Allowed configurations & \Slab \\
Graph hypersurface & Geometry of a Feynman graph & Container of graph integral data & \Slab \\
Feynman integral & Period integral & Observable from geometry $\times$ domain & \Slab \\
Period & Measurable number from geometry & Numerical shadow of structure & \Slab \\
Period pairing $\int_\Gamma\omega$ & Amplitude contribution & Pairing of form and cycle & \Slab \\
Differential form $\omega$ & Integrand / observable density & What is accumulated & \Slab \\
Cycle $\Gamma$ & Integration domain / history class & Where accumulation occurs & \Slab \\
Relative cohomology & Boundary-sensitive observable & Cuts, thresholds, boundaries & \Slab \\
Motive & Pre-numerical functional essence & Hidden source of periods & \Slab/\Hlab \\
Mixed motive & Layered informational object & Multiple weights/depths & \Slab/\Hlab \\
Mixed Tate motive & Polylogarithmic period structure & Tame amplitude period class & \Slab \\
Non-Tate motive & More complex amplitude geometry & Elliptic, Calabi--Yau, modular & \Slab \\
Multiple zeta value & Special period & Constant term in loop expansions & \Slab \\
Polylogarithm $\Li_n(x)$ & Layered propagation function & Iterated contribution & \Slab \\
Multiple polylogarithm & Multi-scale iterated propagation & Nested channel dependence & \Slab \\
Symbol of polylogarithm & First-order function decomposition & Branch-cut/discontinuity data & \Slab \\
Coproduct / coaction & Decomposition of a period & How amplitude splits & \Slab \\
Cobracket $\cobr$ & Antisymmetric primitive split & Factorization/information split & \Hlab \\
Goncharov coproduct & Motivic decomposition operation & Arithmetic anatomy of polylogs & \Slab \\
Goncharov Lie coalgebra & Coalgebra of polylog data & Grammar of primitive amplitude data & \Slab/\Hlab \\
Weight grading & Transcendental complexity & Loop depth/functional complexity & \Slab/\Hlab \\
Depth grading & Iterated-integral nesting & Nested propagation layers & \Slab/\Hlab \\
Primitive element & Indecomposable motivic component & Irreducible contribution & \Hlab \\
Regulator map & Motive to analytic number & Measurement/numerical extraction & \Slab/\Hlab \\
Hodge realization & Analytic-geometric reading & Differential-form representation & \Slab \\
de Rham realization & Form/integrand side & Calculational representation & \Slab \\
Betti realization & Cycle/topological side & Path/integration-domain rep. & \Slab \\
\'Etale realization & Arithmetic/discrete realization & Number-theoretic shadow & \Slab/\Hlab \\
Cosmic Galois group & Symmetry acting on periods & Hidden symmetry of amplitudes & \Hlab/\Plab \\
\end{longtable}

\end{document}
